\section*{LOCALIZACIÓN DE FUENTES CON UN GEIGER MÜLLER}

\textbf{OBJETIVO:} El alumno debe estar preparado en cualquier momento para tener la calma y conocimientos suficientes para poder encontrar y rescatar una fuente radiactiva. Esta práctica debe servir para que los alumnos se familiaricen con los procedimientos de localización de fuentes radiactivas.

\textbf{MATERIAL:}
\begin{enumerate}
    \item Contenedor de Gammagrafía.\\
    (Contiene una fuente radiactiva con la suficiente actividad como para poder tener en sus costados niveles de rapidez de exposición de alrededor de 50 a 100 mR/h).
    \item Equipo portátil detector de Radiaciones tipo Geiger-Müller.
    \item Cronómetro.
\end{enumerate}

\textbf{DESARROLLO:} Se deben formar grupos de trabajo de dos personas para efectuar la búsqueda de la fuente radiactiva. Una vez formados estos grupos se debe esconder el contenedor en un patio o área tal que sea difícil su localización, realizando la operación sin que los alumnos se den cuenta. Una vez escondido éste, se le permite al primer grupo iniciar la búsqueda de la fuente haciendo para ello uso del Detector de Radiaciones. A todas las demás personas que participan en esta práctica, se les permitirá ver cómo sus compañeros encuentran la fuente escondida, analizando los aciertos y errores del grupo de trabajo involucrado en la operación. Se deben tomar los tiempos empleados por cada grupo de trabajo para la localización de la fuente.